\documentclass[11pt]{article}
\usepackage[margin=1in]{geometry}
\usepackage[T1]{fontenc}
\usepackage{lmodern,microtype,amsmath,amssymb,amsthm,booktabs}
\usepackage{xurl}
\usepackage[colorlinks=true,urlcolor=blue,citecolor=blue,linkcolor=blue]{hyperref}
\newtheorem{theorem}{Theorem}
\newtheorem{lemma}[theorem]{Lemma}
\newtheorem{proposition}[theorem]{Proposition}
\newtheorem{remark}[theorem]{Remark}
\newcommand{\R}{\mathbb R}
\newcommand{\supp}{\operatorname{supp}}
\newcommand{\diag}{\operatorname{diag}}
\newcommand{\Hess}{\operatorname{Hess}}
\hypersetup{pdftitle={A sharp forest bound and an improved whisker bound for finite WARM graphs},pdfauthor={Alec Kriebel},pdfsubject={Unrefereed mathematical preprint; finite WARM graph reinforcement},pdfkeywords={WARM, reinforcement, forest, whisker, stability, odd cycle}}
\title{A sharp forest bound and an improved whisker bound\protect\\for finite WARM graphs}
\author{Alec Kriebel\\\small Independent researcher\\
\small ORCID: \href{https://orcid.org/0009-0001-9320-500X}{0009-0001-9320-500X}}
\date{21 October 2026 \quad Version 1.0.0 \quad Preprint}
\begin{document}
\maketitle
\begin{abstract}
Consider power reinforcement with exponent $\alpha>1$ on a finite simple
loopless graph: vertices fire at positive rates and select an incident edge
proportionally to its current count to the power $\alpha$.
We prove that the infinitely reinforced edges almost surely form a forest
for $\alpha>4/3$, for arbitrary positive vertex rates and positive integer
initial counts. The known triangle obstruction shows that the constant
$4/3$ cannot be decreased; its equality case is excluded.
For equal rates we prove that every surviving component has diameter at
most three when $\alpha\ge17/4$.
The main deterministic step excludes an arbitrary weighted odd support
cycle of length $m$ when $\alpha>\sec^2(\pi/(2m))$; the corresponding uniform
cycle threshold and the even-cycle method are known.
We use the published finite WARM convergence and support theorems, and
credit the earlier least-weight and champion arguments.
The forest result answers the forest-bound improvement requested in
Oberwolfach Reports 12/2023; the suggested sharp whisker bound remains open.
AI tools were extensively involved. This is an unrefereed preprint.
\end{abstract}

\section{The problem, model and contribution}
Let $G=(V,E)$ be a finite simple undirected loopless graph. At each
non-isolated vertex $v$ an independent Poisson clock has rate
$\lambda_v>0$. Each edge has an initial count $N_0(e)=n_e\in\{1,2,\ldots\}$.
At a firing of $v$, an incident edge $e$ is selected with probability
\[
 \frac{N_t(e)^\alpha}{\sum_{f\ni v}N_t(f)^\alpha},
\]
and its count is increased by one. Isolated vertices have no effect.
Define
\[
 E_\infty=\{e:N_t(e)\longrightarrow\infty\},\qquad
 E_+=\{e:\liminf_{t\to\infty}N_t(e)/t>0\}.
\]
A whisker is a tree of diameter at most three; an isolated vertex is harmless.

Kleptsyn's contribution to Oberwolfach Reports 12/2023
\cite[printed pp.~653--655]{OWR} asks to improve the sufficient forest
bound $\alpha>2$, allowing unequal rates, and the sufficient equal-rate
whisker bound $\alpha>25$. These bounds come from Holmes--Kleptsyn
\cite{HK}, with the full stochastic bridge supplied by
Hirsch--Holmes--Kleptsyn \cite{HHK}. The report separately suggests that
$4/3$ might be sharp for both conclusions. An improvement request and
this stronger sharpness suggestion are distinct.

\begin{samepage}
\begin{theorem}\label{main}
For every graph and every positive integer initial count vector as above:
\begin{enumerate}
\item If $\alpha>4/3$, then $(V,E_\infty)$ is almost surely a forest,
for arbitrary positive rates.
\item If the rates are equal and $\alpha\ge17/4$, then almost surely every
component of $(V,E_\infty)$ is a whisker.
\end{enumerate}
For each $1<\alpha<4/3$ the equal-rate triangle with unit initial counts
has positive probability that all three edges survive.
Consequently no smaller constant can replace $4/3$ in the first assertion.
No assertion is made at $\alpha=4/3$.
\end{theorem}
\end{samepage}

The new mechanism considered here is the weighted odd-cycle estimate,
which applies to every equilibrium with that support, even with unequal
rates, chords and attachments. The threshold for the \emph{uniform}
equilibrium on a bare odd cycle was already given in
van der Hofstad--Holmes--Kuznetsov--Ruszel
\cite[Proposition~2(iv)]{VHHKR}. The even-cycle cancellation is already
in \cite[Lemma~2]{HHK}, whose proof extends immediately to positive rate
coefficients. The triangle obstruction is also known
\cite[Proposition~2(v), Theorem~4(ii)]{VHHKR}.
The whisker proof sharpens the established least-weight/champion approach.
The earlier \emph{stated} bound is $25$; in fact the displayed constants
in \cite[Section~5.1]{HK}, together with \cite{HHK}, already imply the
weaker improvement $\alpha>24$. A precise derivation is recorded in
Section~\ref{priority}. Thus we do not claim the first possible numerical
improvement of $25$.

\section{The stochastic input and initialization}
Put $\Lambda=\sum_v\lambda_v$, discarding isolated vertices. If $E$ is
empty the conclusions are immediate. Otherwise $\Lambda>0$. Write
\begin{equation}\label{field}
 S_v(a)=\sum_{e\ni v}a_e^\alpha,\qquad
 F_e(a)=-a_e+\sum_{v\ni e}\lambda_v\frac{a_e^\alpha}{S_v(a)}.
\end{equation}
We use counts per unit time, so $\sum_ea_e=\Lambda$ at equilibrium.
The probability normalization $p_v=\lambda_v/\Lambda$, $x=a/\Lambda$
satisfies $F_\lambda(a)=\Lambda F_p(a/\Lambda)$; the Jacobians are identical.

Theorem~1 of \cite[printed p.~3016]{HHK} proves, from unit initial counts
on a finite graph, almost-sure convergence of $(N_t(e)/t)_{e\in E}$ to
an equilibrium whose Jacobian has no positive eigenvalue. It also gives
positive probability of convergence to each strictly linearly stable
equilibrium. Proposition~1 of that paper identifies $E_\infty=E_+$
almost surely under bounded degrees and bounded positive rates; these
hypotheses are automatic for our finite graph.
These published stochastic results are imported, not reproved here.

\begin{lemma}\label{initial}
Every almost-sure tail assertion just described, including the support
conclusions deduced from them, transfers from unit counts to any fixed
positive integer initial count vector.
\end{lemma}
\begin{proof}
In the embedded jump chain, prescribe $n_e-1$ reinforcements of each edge
$e$, in any fixed order. At each step prescribe one endpoint of that edge
to fire. The vertex probability is $\lambda_v/\Lambda>0$, and every
prescribed edge probability is positive because all counts are positive.
The event that this finite prescription occurs has positive probability,
and leaves exactly the count vector $n$.
Conditional on this event, the future chain has precisely the law started
at $n$. The Poisson memoryless property gives the corresponding
continuous-time future law. Failure of an almost-sure tail assertion with
positive probability from $n$ would therefore give failure with positive
probability from unit counts. Finite prefixes and finite time shifts do
not change convergence of counts divided by time, their limit, or which
counts diverge. This proves the assertion.
\end{proof}
This one-way argument does not transfer positive attainment of every
specified stable equilibrium to every initialization. We use attainment
only from the unit state for the known triangle obstruction.
Arbitrary positive noninteger initial offsets, initially zero edges,
loops and infinite graphs are outside the stated model.

At a stochastic limit $a$, each positive-rate vertex has an incident
positive edge: its incident count sum dominates the number of firings
at that vertex, so its limiting incident-weight sum is at least
$\lambda_v$. Thus every $S_v(a)>0$. Restrict to the positive support $H$.
For $\alpha>1$, derivatives of a zero coordinate's $\alpha$th power
vanish at zero. Consequently the full Jacobian is the direct sum of
the positive-support Jacobian and $-I$ on zero edges. This derivative
can equivalently be obtained with the local $C^1$ extension
$x\mapsto(\max\{x,0\})^\alpha$. Hence a positive eigenvalue on $H$
excludes the full boundary equilibrium. All subsequent calculations
are on this positive face.

\section{A weighted cycle instability estimate}
Define
\[
 L(a)=-\sum_ea_e+\frac1\alpha\sum_v\lambda_v\log S_v(a).
\]
Then $F=\diag(a)\nabla L$. At an equilibrium, $\nabla L=0$ and
$DF=\diag(a)\Hess L$, which is similar to the symmetric matrix
$\diag(\sqrt a)\Hess L\diag(\sqrt a)$. Thus a positive Hessian
direction gives a positive real Jacobian eigenvalue, even when other
directions are critical.

Use coordinates $y_e=a_e^\alpha$ and write the same potential as
\[
 \Phi(y)=-\sum_ey_e^{1/\alpha}
       +\frac1\alpha\sum_v\lambda_v\log\sum_{e\ni v}y_e.
\]
At an equilibrium the Hessians are congruent, since the gradient
vanishes. Put
\[
 b_v=\frac{\lambda_v}{S_v},\quad c_v=\frac{\lambda_v}{S_v^2},
 \quad D_e=a_e^{1-2\alpha}.
\]
The equilibrium equation gives
\begin{equation}\label{eq}
 a_e^{1-\alpha}=b_u+b_v\quad(e=uv),
\end{equation}
and direct differentiation gives, for every real edge vector $w$,
\begin{equation}\label{hess}
 \alpha^2\Hess\Phi[w,w]
  =(\alpha-1)\sum_eD_ew_e^2
     -\alpha\sum_vc_v\Big(\sum_{e\ni v}w_e\Big)^2.
\end{equation}
No equal-weight or equal-rate assumption occurs here.

\begin{theorem}\label{cycle}
An equilibrium whose positive support contains an even cycle is
linearly unstable for every $\alpha>1$. If the support contains an odd
cycle of length $m$, it is unstable whenever
$\alpha>\sec^2(\pi/(2m))$.
\end{theorem}
\begin{proof}
Let $e_0,\ldots,e_{m-1}$ be a simple support cycle, and let $v_i$
meet $e_{i-1},e_i$, with indices modulo $m$. Put $A_i=a_{e_i}^\alpha$.
Equation~\eqref{eq} implies
\[
 \sum_iD_{e_i}
  =\sum_ib_{v_i}(A_{i-1}^{-1}+A_i^{-1}).
\]
Moreover
\[
 c_{v_i}=\frac{b_{v_i}}{S_{v_i}}
 \le\frac{b_{v_i}}{A_{i-1}+A_i}
 \le\frac{b_{v_i}}4(A_{i-1}^{-1}+A_i^{-1}),
\]
where the last inequality is $(A_{i-1}-A_i)^2\ge0$ after clearing
positive denominators. Therefore
\begin{equation}\label{cyclebound}
 4\sum_ic_{v_i}\le\sum_iD_{e_i}.
\end{equation}
Extra incident positive edges only increase $S_v$. Set all perturbation
coordinates off the cycle to zero, including any chords.

For even $m$ take $w_{e_i}=(-1)^i$. Every incidence sum in
\eqref{hess} is zero and its diagonal term is strictly positive.
This is the known alternating mechanism of \cite[Lemma~2]{HHK}, with
positive vertex coefficients.

For odd $m$ set $\theta=\pi-\pi/m$ and use the two real vectors
$u_{e_i}=\cos(i\theta)$ and $w_{e_i}=\sin(i\theta)$.
Since $m\theta=(m-1)\pi$ is a multiple of $2\pi$, they are cyclically
consistent. At each edge $u_{e_i}^2+w_{e_i}^2=1$, while at each cycle
vertex
\[
 (u_{e_{i-1}}+u_{e_i})^2+(w_{e_{i-1}}+w_{e_i})^2
 =4\sin^2(\pi/(2m)).
\]
Adding \eqref{hess} for the two vectors and using
\eqref{cyclebound} yields
\[
 \alpha^2\big(\Hess\Phi[u,u]+\Hess\Phi[w,w]\big)
 \ge\big[\alpha\cos^2(\pi/(2m))-1\big]\sum_iD_{e_i}>0.
\]
At least one vector is a positive Hessian direction, proving instability.
\end{proof}
Every finite nonforest has a simple cycle, and for odd $m\ge3$
we have $\sec^2(\pi/(2m))\le4/3$. The stochastic input, boundary
argument and Lemma~\ref{initial} prove Theorem~\ref{main}(1).
There is no spurious positive radial eigenvalue: $\sum_eF_e=\Lambda-\sum_ea_e$
implies $\mathbf1^TDF=-\mathbf1^T$, so positive-eigenvalue vectors lie
in the tangent space $\sum_eh_e=0$.

For the known sharpness obstruction take a triangle with all rates
and limiting weights equal to one. Its Jacobian has diagonal
$\alpha/2-1$ and off-diagonal $-\alpha/4$, giving eigenvalues
$-1$ and $3\alpha/4-1$ twice. It is strictly stable for
$1<\alpha<4/3$, and the published attainment theorem gives positive
probability from unit counts. All three edges survive. At $4/3$ the
two transverse eigenvalues vanish, and neither this attainment clause
nor our strict instability argument decides the stochastic outcome.

\section{Equal-rate whiskers}
Normalize the common rate to one by time scaling. For $e=uv$ write
\[
 q=\frac{a_e^\alpha}{S_u},\qquad r=\frac{a_e^\alpha}{S_v}.
\]
At equilibrium $a_e=q+r$, where $0<q,r\le1$, hence $a_e\le2$.
At an equilibrium with no positive eigenvalue, the symmetric Jacobian
representative is negative semidefinite. Its diagonal entries give
\begin{equation}\label{diagonal}
 q^2+r^2\ge(1-1/\alpha)(q+r).
\end{equation}
For example, direct differentiation gives
$J_{ee}=\alpha-1-\alpha(q^2+r^2)/a_e$, proving the condition.
It is necessary, not sufficient, for stability.

\begin{lemma}[Least-weight argument of \cite{HK}]\label{least}
For $\alpha>2$, every positive edge at such an equilibrium has weight
at least one.
\end{lemma}
\begin{proof}
Choose an edge of least positive weight in one component. If neither
endpoint is a leaf of the support, each has a competing edge of at least
the same weight. Then $q,r\le1/2$, so
$q^2+r^2\le(q+r)/2$, contradicting \eqref{diagonal}.
Thus one endpoint is a leaf and contributes one to the weight.
All other weights in the component are no smaller.
\end{proof}

Call an edge internal if neither endpoint is a support leaf. Suppose
an internal edge has weight $a\in[1,3/2]$. Each endpoint has a
competitor of weight at least one, so
\begin{equation}\label{cap}
 q,r\le s_\alpha(a):=\frac{a^\alpha}{1+a^\alpha},\qquad q+r=a.
\end{equation}
For a feasible split $s\ge a/2$, convexity on $[a-s,s]$ gives
\begin{equation}\label{f}
 \frac{q^2+r^2}a\le f(a,s_\alpha(a)),\qquad
 f(a,s)=\frac{s^2+(a-s)^2}a.
\end{equation}
The function $f$ increases in $s$ on $s\ge a/2$, and $s_\alpha(a)$
is nondecreasing in $\alpha$ for $a\ge1$.

\subsection{The closed interval $17/4\le\alpha\le7$}
For $3\le B\le7$ and $1\le a\le3/2$, one has
\begin{equation}\label{z}
 \frac12\le\frac{s_B(a)}a\le\frac23.
\end{equation}
Indeed the lower bound is equivalent to $a^{B-1}(2-a)\ge1$, which
follows from
\[
 a^2(2-a)-1=(a-1)(1+a-a^2)\ge0.
\]
For the upper bound use $s_B\le s_7$ and
$R(a)=2a^7-3a^6+2>0$. Its derivative is $2a^5(7a-9)$,
so its minimum for $a\ge0$ occurs at $9/7$, with
$R(9/7)=52763/823543>0$.

Put $z=s_B(a)/a$. Direct differentiation and \eqref{z} give
\[
 \frac{d}{da}f(a,s_B(a))
 =1-2z^2+2(2z-1)s_B'(a)\ge1/9>0.
\]
If $A\le\alpha\le B$, it follows that
\begin{equation}\label{endpoint}
 \frac{q^2+r^2}a\le f(a,s_B(a))\le f(3/2,s_B(3/2)).
\end{equation}
For $T\ge(3/2)^B$, put $S=T/(1+T)$. The table gives exact
positive margins $\delta=(1-1/A)-f(3/2,S)$:
\begin{center}
\small
\begin{tabular}{ccccc}
\toprule
$[A,B]$ & $T$ & $S$ & $f(3/2,S)$ & $\delta$\\
\midrule
$[17/4,13/3]$ & $29/5$ & $29/34$ & $1325/1734$ & $1/1734$\\
$[13/3,9/2]$ & $25/4$ & $25/29$ & $3869/5046$ & $163/65598$\\
$[9/2,5]$ & $243/32$ & $243/275$ & $117039/151250$ & $5399/1361250$\\
$[5,6]$ & $729/64$ & $729/793$ & $991335/1257698$ & $74117/6288490$\\
$[6,7]$ & $2187/128$ & $2187/2315$ & $8580639/10718450$ & $527104/16077675$\\
\bottomrule
\end{tabular}
\end{center}
The two nonintegral power caps follow exactly from
\[
 (29/5)^3-(3/2)^{13}=\frac{504313}{1024000}>0,
 \qquad
 (25/4)^2-(3/2)^9=\frac{317}{512}>0.
\]
The other caps are equalities. All $S\ge3/4$, so the second-argument
monotonicity is applicable. The five closed intervals cover $[17/4,7]$,
and \eqref{endpoint} is strictly less than
$1-1/A\le1-1/\alpha$, contradicting \eqref{diagonal}.
Thus no internal edge has weight at most $3/2$ in this exponent interval.

\subsection{The unbounded tail $\alpha\ge7$}
Write $a=1+r_0$, with $0\le r_0\le1/2$. Since $q,r\le1$
and $q+r=a$, convexity gives $q^2+r^2\le1+r_0^2$.
Condition~\eqref{diagonal} implies
\begin{equation}\label{quadratic}
 h(r_0):=\alpha r_0^2-(\alpha-1)r_0+1\ge0.
\end{equation}
For $\alpha\ge7$, $h(0)>0$ and $h(1/2)=(6-\alpha)/4<0$.
Thus \eqref{quadratic} forces $r_0$ below the smaller root of $h$.
Set $C=3-\sqrt2$. We have $1<C<8/5$, $C/\alpha<1/2$ and
$C^2+C=7(C-1)$, whence
\[
 h(C/\alpha)=-(C-1)(1-7/\alpha)\le0.
\]
Therefore $r_0\le C/\alpha$ and
\[
 a^\alpha\le\exp(\alpha r_0)\le e^C<6\le\alpha-1.
\]
For completeness the strict bound is elementary: $\sqrt2>7/5$ gives
$C<8/5$. The exponential series at $8/5$ through degree four sums
to $9067/1875$; its tail from degree five has successive ratios at
most $4/15$, so is at most $4096/34375$. Their sum is
$510973/103125<6$.
By \eqref{cap} both $q,r<(\alpha-1)/\alpha$. But
\eqref{diagonal} says that their weighted mean
$(q^2+r^2)/(q+r)$ is at least $(\alpha-1)/\alpha$, a contradiction.
This proves the same internal-edge exclusion for every $\alpha\ge7$.

\subsection{The matching conclusion}
Every internal edge consequently has weight greater than $3/2$.
For such an edge $e=uv$, its probability at $v$ is greater than $1/2$,
since its other endpoint contributes at most one to $a_e=q+r$.
Two internal edges cannot meet at a vertex, where their probabilities
would sum to more than one. Hence internal edges form a matching.
Every cycle has adjacent internal edges and is excluded directly.
A tree of diameter at least four also has adjacent internal edges:
take the two middle edges of a four-edge path. Every support component
therefore has diameter at most three. This proof of the whisker bound
does not depend on the new odd-cycle theorem. Apply the published
stochastic input and Lemma~\ref{initial} to finish Theorem~\ref{main}(2).

\section{Priority, attribution and precise limits}\label{priority}
The main comparison is with the arbitrary-rate forest threshold in
\cite[Theorem~3 and Corollary~1]{HK}, the equal-rate whisker theorem
\cite[Theorem~2(a)]{HK}, the uniform odd-cycle calculation in
\cite[Proposition~2(iv)]{VHHKR}, and the even-cycle proof in
\cite[Lemma~2]{HHK}. We make no new claim for the stochastic convergence
or support-identification theorems, the triangle obstruction, or the
least-weight/champion mechanism. A bounded current literature audit
recovered no earlier theorem giving the arbitrary-equilibrium weighted
odd-cycle estimate or either substantive bound in Theorem~\ref{main}.
This is not a certificate of exhaustive priority or worldwide firstness.
Repository creation dates do not certify historical public visibility.

The earlier published argument already has a small numerical improvement
of its stated $25$. Here is an explicit bridge, to avoid understating it.
In \cite[Section~5.1, Lemmas~5--7]{HK} put $c=9/4$.
For $\alpha>81/4$, its diagonal/least-weight argument gives
\[
 a_e\in[1,1+c/\alpha)\ \cup\ (2-2c/\alpha,2].
\]
An edge in the first interval must have a leaf endpoint once
$\alpha>c(e^c+1)$: otherwise its champion probability is at most
$e^c/(e^c+1)$ instead of greater than $1-c/\alpha$.
The exact bound
\[
 e^{9/4}\le\sum_{k=0}^{12}\frac{(9/4)^k}{k!}
 +\frac{(9/4)^{13}/13!}{1-(9/4)/14}
 =\frac{191714244173193067}{20206532637491200}<\frac{19}2
\]
gives $c(e^c+1)<189/8<24$. For $\alpha>24$ the second interval has
lower endpoint greater than $29/16>3/2$, so these edges form a matching
and the other edges have leaf endpoints. The published convergence
and support bridge then yields whiskers. Thus the purely literal request
for some improvement of $25$ has this earlier consequence; our
contribution is the stronger $17/4$ bound, alongside the forest result.

The interval $4/3<\alpha<17/4$ in the suggested sharp equal-rate
whisker conclusion is unresolved by this paper. The stochastic equality
case $\alpha=4/3$ is also unresolved here. Unequal-rate whisker conclusions
are false in general \cite[Example~1]{HK}. We assert neither an infinite-graph
extension nor results for loops or noninteger initial offsets.
The finite simple graph model and all quantifiers appear explicitly above.

\section{Reproducibility and research disclosure}
The proof, including every rational margin and the unbounded tail, is
contained in this paper and does not depend on numerical searches.
The supporting package preserves the eight original PR artifacts,
separate mathematical and priority audits, exact finite controls and
instructions for reproduction in a disposable copy. Historical source
wrappers label the work partial by requiring the sharper suggested
whisker bound; that label is explicitly superseded here for the literal
improvement requests. The added integer-count conditioning argument and
prior-method credits are identified, and original inputs remain unchanged.
Finite rational checks and rejected deliberately false alternatives
corroborate the derivations; hashes establish custody rather than proof.
No exploratory finite tree search is used as universal evidence.
Third-party primary texts and private operational captures are not
redistributed. Standard background results are cited at their interfaces.

AI tools were extensively involved in mathematical research, proof
development, symbolic and computational checks, adversarial auditing,
literature investigation and manuscript preparation. Independent internal
AI reviews are not independent human peer review. This manuscript is a
preprint and has not undergone formal independent human peer review,
journal acceptance or proof-assistant certification.

\begin{thebibliography}{9}
\bibitem{OWR} V.~Kleptsyn, Graph-based interacting P\'olya urns,
in
\emph{MATRIX--MFO Tandem Workshop: Stochastic Reinforcement Processes and
Graphs}, Oberwolfach Reports 12/2023, printed pp.~653--655,
Open Problems~1--2 on p.~655.
\href{https://doi.org/10.4171/owr/2023/12}{doi:10.4171/owr/2023/12}.
\url{https://ems.press/content/serial-article-files/47008}.
\bibitem{VHHKR} R.~van der Hofstad, M.~Holmes, A.~Kuznetsov and W.~Ruszel,
Strongly reinforced P\'olya urns with graph-based competition,
\emph{Ann. Appl. Probab.}~26 (2016), 2494--2539.
\href{https://doi.org/10.1214/16-AAP1153}{doi:10.1214/16-AAP1153}.
\bibitem{HK} M.~Holmes and V.~Kleptsyn,
Proof of the WARM whisker conjecture for neuronal connections,
\emph{Chaos}~27 (2017), 043104.
\href{https://doi.org/10.1063/1.4978683}{doi:10.1063/1.4978683}.
Author manuscript:
\url{https://researchers.ms.unimelb.edu.au/~mholmes1@unimelb/Reinforced-graphs_final.pdf}.
\bibitem{HHK} C.~Hirsch, M.~Holmes and V.~Kleptsyn,
Infinite WARM graphs III: strong reinforcement regime,
\emph{Nonlinearity}~36 (2023), 3013--3042.
\href{https://doi.org/10.1088/1361-6544/acc9a0}{doi:10.1088/1361-6544/acc9a0}.
\end{thebibliography}
\end{document}
